\SourcePage{597}{600}
\section{Isotope shift}

Properties specific to the nucleus---its finite mass, dimensions, and spin---
distinguish it from a fixed point centre of a Coulomb field and have a
definite effect on the electronic energy levels of an atom.

One such effect is the so-called \emph{isotope shift} of the levels: the change
in a level's energy on passing from one isotope of an element to another. What
is of physical interest is, of course, not the change in a single level but the
change in the separation of two levels, which is observed as a spectral line.
Consequently, it is not the energy of the entire electronic shell that must be
considered, but only the part of it associated with the electron taking part in
the spectral transition in question.

For light atoms the chief origin of the isotope shift is the finiteness of the
nuclear mass. Allowance for the motion of the nucleus adds to the atomic
Hamiltonian the term
\[
 \frac1{2M}\left(\sum_i\hat{\mathbf p}_i\right)^2,
\]
where $M$ is the mass of the nucleus and $\mathbf p_i$ are the electron
momenta\footnote{In the atom's centre-of-mass frame, the momenta of the nucleus
and the electrons have zero sum:
$\mathbf p_{\mathrm{nuc}}+\sum_i\mathbf p_i=0$. Their total kinetic energy is
therefore
\[
 \frac{\mathbf p_{\mathrm{nuc}}^2}{2M}+
 \frac1{2m}\sum_i\mathbf p_i^2
 =\frac1{2M}\left(\sum_i\mathbf p_i\right)^2
 +\frac1{2m}\sum_i\mathbf p_i^2.
\]
}. The isotope shift due to this effect is thus the mean value
\begin{equation}
 \frac12\left(\frac1{M_1}-\frac1{M_2}\right)
 \overline{\left(\sum_i\mathbf p_i\right)^2},
 \tag{120.1}
\end{equation}
evaluated with the wave function of the atomic state concerned; $M_1$ and
$M_2$ are the nuclear masses of the two isotopes.

\SourcePage{598}{601}
For heavy atoms the principal contribution to the isotope shift comes from the
finite extent of the nucleus. In practice, this effect is appreciable only for
levels of an outer electron in an $s$-state. Indeed, unlike the wave functions
of states with $l\ne0$, the $s$-state wave function does not vanish as
$r\to0$, so the probability of finding the electron within the ``volume of
the nucleus'' is comparatively large. We shall calculate the isotope shift in
this case\unskip
\footnote{The calculation below does not include relativistic effects. The omitted effects concern the electron's motion near the nucleus. It is valid under the condition $Ze^2/(\hbar c)\ll1$.}
.

Let $\varphi(r)$ denote the actual electrostatic potential of the nuclear
field, as distinct from the Coulomb potential $Ze/r$ of a point charge $Ze$.
The resulting change in the electron's energy relative to its value in the
pure Coulomb field $Ze/r$ is
\begin{equation}
 \Delta E=-e\int\left(\varphi-\frac{Ze}{r}\right)\psi^2(r)\,dV,
 \tag{120.2}
\end{equation}
where $\psi(r)$ is the electron wave function (for an $s$-state it is real and
spherically symmetric). Although the integral is formally over all space, the
difference $\varphi-Ze/r$ in its integrand is nonzero only inside the nuclear
volume.

On the other hand, the wave function of an $s$-state tends to a constant as
$r\to0$ (see \S~32), and this limiting value is already reached, to sufficient
accuracy, outside the nucleus. We may therefore take $\psi^2$ outside the
integral, replacing $\psi(r)$ by its value at $r=0$ as calculated for the
Coulomb field of a point charge.

To transform the remaining integral, use the identity $\Delta r^2=6$ and
rewrite (120.2) as
\[
 \Delta E=-\frac16e\psi^2(0)\int
 \left(\varphi-\frac{Ze}{r}\right)\Delta r^2\,dV
 =-\frac16e\psi^2(0)\int r^2\Delta
 \left(\varphi-\frac{Ze}{r}\right)dV.
\]
In this transformation of the volume integral, the surface integral at
infinity has been set equal to zero. Now
$\Delta(1/r)=-4\pi\delta(\mathbf r)$, while
$r^2\delta(\mathbf r)=0$ for all $r$. The electrostatic Poisson equation gives
$\Delta\varphi=-4\pi\rho$, where here $\rho$ is the density of electric charge
in the nucleus. It follows finally that
\begin{equation}
 \Delta E=\frac{2\pi}{3}\psi^2(0)Ze^2\overline{r^2},
 \tag{120.3}
\end{equation}

\SourcePage{599}{602}
where
\[
 \overline{r^2}=\frac1{Ze}\int\rho r^2\,dV
\]
is the nuclear proton mean-square radius. (For a uniform distribution of
protons in the nucleus it would be $\overline{r^2}=3R^2/5$, with $R$ the
geometrical nuclear radius.) The isotope shift of a level is determined by the
difference between (120.3) for the two isotopes.

In \S~71 the magnitude of $\psi(0)$ was estimated and found to depend on the
atomic number (assumed large) as $\sqrt Z$. The splitting in (120.3) is
therefore proportional to $R^2Z^2$.
